% ===================================================================
% appendices/B_code.tex -- Source code (writer C2; filled by W4). ASCII only.
% Listings are the ASCII-sanitised copies in code/ made by tools/make_appendix_assets.py.
% ===================================================================
\chapter{Source code}\label{ch:app-code}

This appendix reproduces every Python script of the package that generates decks, launches the solver, extracts metrics
or produces analysis numbers quoted in the report, so that each number can be re-derived from code rather than from
prose. The listings are copies of the package files, made ASCII-safe for typesetting by
\file{tools/make_appendix_assets.py}; the package files themselves (\file{PKG/scripts/} and
\file{PKG/analysis_2026-09-25/scratch/}) are authoritative. The commands used to run them are collected in
\cref{ch:app-repro}.

\section{Organisation of the code}\label{sec:code-overview}

The code falls into three groups. The first group consists of the package scripts in \file{scripts/}
(\cref{sec:code-package}): the material laws, the deck generator, the runner and campaign driver that launch the real
solver, the common metric extractor, and the table and plot writers. Only \file{run_atlas.py}, \file{run_campaign.py}
and \file{sensitivity_analysis.py} launch ATLAS; all other scripts read existing files. The second group consists of the
analysis scripts of 2026-09-25 in \file{analysis_2026-09-25/scratch/<S1..S6, REVIEW, SYNTHESIS>/}
(\cref{sec:code-analysis}), which contain the specialist analyses. None of them launches ATLAS or modifies a package
file; they read run outputs and the measured data, and several import \file{scripts/extract_metrics.py} unchanged, so
that the same extraction is used everywhere. In the listings below the file name is prefixed with the analysis folder
(e.g.\ \file{S1_s1_analytic.py} is \file{scratch/S1/s1_analytic.py}). The third group consists of the report tools in
\file{REPORT/tools/} (\cref{sec:code-collect}), which build the figures, the digest and the appendix assets of this
report; they are described but not listed.

Results of surrogate models (1-D Poisson, Schr\"odinger-Poisson, charge-sheet models) are always labelled \evSURR{} in the
report. They are used for decomposition and hypothesis generation and never as a substitute for an ATLAS curve.

% -------------------------------------------------------------------
\section{Package scripts}\label{sec:code-package}

\subsection{\texttt{material\_laws.py}}
The script implements \texttt{IWO\_material(t, xW)}, the thickness-aware material laws of about 2\,at.\,\% W:\InO{} used
by V1. It fits the thickness laws to the literature and DFT evidence (confinement shift, effective mass, tail density,
background donors, roughness factor) and evaluates them at a requested thickness; no simulator is run. The deck generator
imports it, so that every deck is produced from the same laws (\cref{ch:derivations-es,ch:derivations-tr}).
\lstinputlisting[language=Python]{code/scripts/material_laws.py}

\subsection{\texttt{build\_iwo\_decks.py}}
The script generates the ATLAS decks from one material model, \texttt{IWO\_material(t, xW)}, for the requested
thicknesses. It accepts dotted-path overrides of the material model (\texttt{--override}), the structure form
(\texttt{--structure full|reduced}) and the sweep form (\texttt{--sweep stepped|pointwise}), and it writes the provenance
header of every deck (\cref{sec:num-deck}). All decks of \cref{ch:app-decks} were produced with it.
\lstinputlisting[language=Python]{code/scripts/build_iwo_decks.py}

\subsection{\texttt{run\_atlas.py}}
The script launches one real ATLAS run of a generated deck through DeckBuild in batch mode. It validates the terminal
currents (status, KCL, drain sign, native zeros), scores the curve against the measurement with the common extractor and
appends one row to \file{results/RUN_INDEX.csv}; a Python model is never substituted for the solver. Options select the
thickness, label, overrides, sweep segments, structure and timeout. The option \texttt{--rescore} repeats validation and
scoring from an existing run folder without a launch, and \texttt{--output-vgs} adds the $\Id$--$\Vd$ families of the
prediction runs (\cref{sec:num-validation,sec:num-campaign-runner}).
\lstinputlisting[language=Python]{code/scripts/run_atlas.py}

\subsection{\texttt{run\_campaign.py}}
The script runs a sequential ATLAS campaign from a JSON plan (the 2026-09-25 hypothesis, convergence, sensitivity and
prediction study), e.g.\ \texttt{python scripts/run\_campaign.py config/campaign\_A\_6p3\_hypotheses.json}. Each plan
entry is launched through \file{run_atlas.py}, one at a time, and the metrics of each run are collected in
\file{results/campaigns/<name>.json}.
\lstinputlisting[language=Python]{code/scripts/run_campaign.py}

\subsection{\texttt{extract\_metrics.py}}
This is the single metric extractor, applied identically to measured and simulated $\Id$--$\Vg$ curves; no simulator is
called. Its docstring states the conventions used for $\Vthcc$, $\Vthlin$, the swings, $\gmmax$, $\muFE$ and $\Ion$,
because the target paper does not define its extraction (\cref{sec:der-metrics}). Called with \texttt{--thickness}, it
prints the metrics of a measured curve; with \texttt{--sim}, it applies the simulation conventions.
\lstinputlisting[language=Python]{code/scripts/extract_metrics.py}

\subsection{\texttt{compare\_experiment\_simulation.py}}
The script builds \file{tables/EXPERIMENTAL_VS_SIMULATION.{csv,md}} and the overlay and residual plots of each thickness
from the best V1 runs listed in \file{results/BEST_RUNS_V1.json}, with the same extraction for measured and simulated
curves. No simulator is called.
\lstinputlisting[language=Python]{code/scripts/compare_experiment_simulation.py}

\subsection{\texttt{make\_provenance\_tables.py}}
The script writes \file{tables/MATERIAL_PARAMETERS_USED.{csv,md}}, \file{tables/PARAMETER_PROVENANCE.csv} and
\file{docs/PARAMETER_PROVENANCE.md} from \file{config/iwo_material_model.yaml}, optionally with the calibrated overrides
recorded in \file{results/BEST_RUNS_V1.json}. These tables form the basis of \cref{tab:provenance-master}.
\lstinputlisting[language=Python]{code/scripts/make_provenance_tables.py}

\subsection{\texttt{sensitivity\_analysis.py}}
The script performs a one-at-a-time sensitivity analysis of the V1 model in the real solver; each case is one ATLAS
launch relative to the best run of a thickness (e.g.\ \texttt{--thickness 2 --cases chi\_m0p1 nd\_x2 dit\_x3}). It
produced the sensitivity runs \run{0020}--\run{0027} (\run{0027} malformed and retracted, \cref{sec:cal-param-control});
the chain was stopped by the user when the launch budget was limited.
\lstinputlisting[language=Python]{code/scripts/sensitivity_analysis.py}

\subsection{\texttt{sensitivity\_from\_stage\_runs.py}}
The script derives sensitivities from the real calibration-stage runs already in \file{results/} (pairs of stage runs
that differ in one parameter), without calling the simulator. It was written after the one-at-a-time chain had been
stopped, in order to extract the remaining sensitivities from runs that already existed.
\lstinputlisting[language=Python]{code/scripts/sensitivity_from_stage_runs.py}

\subsection{\texttt{plot\_results.py}}
The script produces plots from the best V1 runs (\file{results/BEST_RUNS_V1.json}) and the thickness laws without calling
the simulator. For each thickness it draws the overlays on logarithmic and linear axes, the residual, the band diagrams
at equilibrium, threshold and on-state, and the electron density versus gate voltage. Several of these PNG files are
reproduced in \cref{ch:derivations-es}.
\lstinputlisting[language=Python]{code/scripts/plot_results.py}

\subsection{\texttt{plot\_sensitivity.py}}
The script writes \file{plots/sensitivity_overview.png}, which shows the no-confinement counterfactual overlays at 2 and
\SI{13.2}{nm} and the one-at-a-time metric changes at \SI{2}{nm} from \file{tables/SENSITIVITY_OAT.csv}. No simulator is
called.
\lstinputlisting[language=Python]{code/scripts/plot_sensitivity.py}

\subsection{\texttt{qe\_build\_inputs.py}}
The script builds an \emph{unexecuted} Quantum ESPRESSO workflow \citep[p.~1]{Giannozzi2009} for a matrix of W content
and slab thickness, writing \texttt{pw.x} input files and a README. Quantum ESPRESSO is not installed on the project
machine, so no DFT number of this report comes from this script; it documents how the DFT proxies of the thickness laws
could be replaced by IWO-specific calculations.
\lstinputlisting[language=Python]{code/scripts/qe_build_inputs.py}

\subsection{\texttt{export\_report\_pdf.py}}
The script renders the earlier Markdown report \file{docs/FULL_REPORT.md} to HTML with print CSS and embedded key
figures, and to PDF with the installed Microsoft Edge in headless mode. It changes no report text and calls no
simulator; it is not part of the present \LaTeX{} report.
\lstinputlisting[language=Python]{code/scripts/export_report_pdf.py}

% -------------------------------------------------------------------
\section{Analysis scripts of 2026-09-25}\label{sec:code-analysis}

\subsection{\texttt{S1\_s1\_analytic.py}}
The script contains the order-of-magnitude analytic electrostatics of specialist S1, with constants taken from the
evidence brief and the material model (\evTHEORY). It provides the dimensional estimates against which the ATLAS-based
S1 tables are checked.
\lstinputlisting[language=Python]{code/analysis/S1_s1_analytic.py}

\subsection{\texttt{S1\_s1\_atlas\_tables.py}}
The script builds the S1 electrostatics tables only from existing ATLAS outputs and the measured data
(\file{execution.json}, \file{comparison.csv} and the conduction-band, quasi-Fermi-level and electron-density profiles of
the runs). No ATLAS launch is made.
\lstinputlisting[language=Python]{code/analysis/S1_s1_atlas_tables.py}

\subsection{\texttt{S1\_s1\_poisson1d.py}}
The script has two parts: (A) post-processing of the ATLAS depth profiles into a trapped-charge budget at the cut
$x = 12\,\mu$m, and (B) a 1-D Poisson plus gradual-channel (Pao-Sah) surrogate of the V1 model, used only for
decomposition and what-if questions (lever arm of a back-surface charge, neutral-layer onset, no-confinement symmetric
test). The surrogate is checked against ATLAS before use; it underlies the threshold budget of \cref{sec:der-vth-budget}
(\evSURR).
\lstinputlisting[language=Python]{code/analysis/S1_s1_poisson1d.py}

\subsection{\texttt{S1\_s1\_shape\_test.py}}
The script addresses S1 task 3, namely which kind of gate-voltage-dependent charge could reproduce the \SI{6.3}{nm}
on-state shape ($\Vthlin$, $\gmmax$) while keeping $\Vthcc$, $\SScc$ and $\Ion$. An extra acceptor Gaussian band (sheet
density, depth below $\Ec$, width) is added to the \run{0015} configuration in the ATLAS-checked 1-D surrogate. The
result is hypothesis-generating only (\evSURR) and motivated the pre-registered test C1 (\run{0052}).
\lstinputlisting[language=Python]{code/analysis/S1_s1_shape_test.py}

\subsection{\texttt{S2\_s2\_part1\_curves.py}}
This is part 1 of the transport analysis of specialist S2. It extracts mobility curves from the measured and simulated
transfer curves and fits the thickness laws of mobility and on-current, including the leave-one-out comparison of power
law, $A + B/t$, exponential and step models (\cref{sec:der-powerlaw}).
\lstinputlisting[language=Python]{code/analysis/S2_s2_part1_curves.py}

\subsection{\texttt{S2\_s2\_part2\_mtr\_temperature.py}}
Part 2 of S2 treats the temperature dependence of the drain current in the multiple-trapping-and-release picture, which
is used to frame the two temperature variants of the predictions (\cref{sec:der-activation,sec:par-temperature}).
\lstinputlisting[language=Python]{code/analysis/S2_s2_part2_mtr_temperature.py}

\subsection{\texttt{S2\_s2\_part3\_repartition.py}}
Part 3 of S2 performs the shape-consistent re-partition of the field-effect mobility into band mobility and
trap-limited factor for each film, which led to the pre-registered B0 test (\run{0037}; \cref{ch:6p3}).
\lstinputlisting[language=Python]{code/analysis/S2_s2_part3_repartition.py}

\subsection{\texttt{S2\_s2\_part4\_musat\_location.py}}
Part 4 of S2 determines the gate voltage at which the saturation-formula mobility peaks. This is relevant to the
interpretation of the mobility values of the paper as saturation-formula peaks near threshold (\cref{sec:der-mufe}).
\lstinputlisting[language=Python]{code/analysis/S2_s2_part4_musat_location.py}

\subsection{\texttt{S3\_s3\_e1\_magnitude.py}}
The script addresses S3 task 1, the magnitude of electron confinement in 2, 6.3 and \SI{13.2}{nm} IWO films, from pure
effective-mass estimates (finite-difference eigenvalues) without fitting to device data (\evTHEORY;
\cref{sec:der-confinement}).
\lstinputlisting[language=Python]{code/analysis/S3_s3_e1_magnitude.py}

\subsection{\texttt{S3\_s3\_necessity\_curves.py}}
The script addresses S3 task 3, whether confinement is required by the transfer data. It uses only existing ATLAS
outputs (\file{idvg.dat}, native A/$\mu$m) and the measured curves, and it starts with a rigid-shift test of the
horizontal distance between the confinement-on and confinement-off curves versus current level (\run{0012}/\run{0016},
\run{0013}/\run{0017}).
\lstinputlisting[language=Python]{code/analysis/S3_s3_necessity_curves.py}

\subsection{\texttt{S3\_s3\_sp1d.py}}
The script addresses S3 task 2 with a minimal self-consistent 1-D Schr\"odinger-Poisson model, compared with a classical
3-D Fermi-Dirac model of the gate stack at the channel centre at $\Vd = 0$ (equilibrium with the source/drain Fermi
level). It produces the onset shifts of \cref{sec:der-confinement} (\evSURR).
\lstinputlisting[language=Python]{code/analysis/S3_s3_sp1d.py}

\subsection{\texttt{S3\_s3\_sp1d\_series.py}}
The script addresses S3 task 5, the same trap-free comparison of Schr\"odinger-Poisson and classical models for a denser
thickness series (3, 4 and \SI{5}{nm}); it imports \file{s3_sp1d.py}.
\lstinputlisting[language=Python]{code/analysis/S3_s3_sp1d_series.py}

\subsection{\texttt{S4\_s4\_traps.py}}
The script contains the defect and trap analysis of specialist S4 and is read-only with respect to the runs. Its parts
include the thickness law of the tail density (power law versus a surface-plus-bulk decomposition, with the numbers of
\citet[pp.~7, 9]{Januar2026} and \citet[pp.~173507-3, 173507-6]{Kim2024}); the trap-spectroscopy figure of
\cref{ch:derivations-es} (\file{scratch/S4/s4_trap_spectroscopy.png}) belongs to the same S4 analysis.
\lstinputlisting[language=Python]{code/analysis/S4_s4_traps.py}

\subsection{\texttt{S4\_s4\_part2.py}}
Part 2 of S4 computes the integrated trap-charge excess at \SI{6.3}{nm} by energy window (measured versus \run{0015},
with the same inversion), the channel current at $\Vg = 0$ compared with the off-floor, and tables of the ratio of
measured to ATLAS-inverted trap densities for each film.
\lstinputlisting[language=Python]{code/analysis/S4_s4_part2.py}

\subsection{\texttt{S4\_s4\_part3\_63shape.py}}
Part 3 of S4 asks which trap redistribution reproduces the \SI{6.3}{nm} on-state shape. It uses a zero-dimensional
charge-sheet surrogate (uniform film potential, Fermi-Dirac free electrons, thermal trap occupancy, Pao-Sah-type integral
at $\Vd = \SI{0.7}{V}$), with each variant rigidly aligned to the measured $\Vthcc = \SI{0.644}{V}$. It proposed the
near-$\Ec$ acceptor band tested as C1 (\run{0052}) (\evSURR).
\lstinputlisting[language=Python]{code/analysis/S4_s4_part3_63shape.py}

\subsection{\texttt{S5\_s5\_yfunction.py}}
The script performs the series-resistance, Y-function and contact analysis of the measured transfer curves and, as a
method check, of ideal-contact ATLAS runs ($\Rsd = 0$ by construction, constant $\muband$). The inputs are the measured
data and the \file{comparison.csv} files of the runs; no parameter is invented (\cref{sec:der-yfunction}).
\lstinputlisting[language=Python]{code/analysis/S5_s5_yfunction.py}

\subsection{\texttt{S5\_s5\_checks.py}}
The S5 checks comprise (1) the ideal gradual-channel baseline of the H-function exponent in the same windows (at
$\Vd = \SI{0.7}{V}$ the condition $\Vd \ll V_{\mathrm{ov}}$ does not hold), (2) the forward effect of a series resistance
on the ideal curve and on a measured-like curve, and (3) the determination of which mobility definition reproduces the
values 5.1 and \SI{27.4}{cm^2/(V.s)} of the paper from the workbook curves (the provenance test of
\cref{sec:inp-metrics-measured}).
\lstinputlisting[language=Python]{code/analysis/S5_s5_checks.py}

\subsection{\texttt{S5\_s5\_plot.py}}
The script draws the Y-function and H-function (threshold-free exponent) figure for the measured curves and the
ideal-contact ATLAS runs, reproduced as the \texttt{yfunction} figure.
\lstinputlisting[language=Python]{code/analysis/S5_s5_plot.py}

\subsection{\texttt{S6\_s6\_lib.py}}
The script contains the common helpers of specialist S6 (read-only): metrics from fixed-current crossings with both
log-linear and PCHIP interpolation on $\log_{10}\Id$ (their difference is reported as interpolation uncertainty),
fixed-current swings and on-state shape markers. The package extractor \file{scripts/extract_metrics.py} is imported
unchanged.
\lstinputlisting[language=Python]{code/analysis/S6_s6_lib.py}

\subsection{\texttt{S6\_s6\_part1\_numerics.py}}
Part 1 of S6 establishes the numerical closure from existing run outputs: linearity in $\muband$, DOS offsets versus gate
voltage and thickness, the 6.3 and \SI{13.2}{nm} mesh checks, KCL, native zeros, the fragility of $\SSmin$, and the
uncertainty due to interpolation and to the transconductance grid. It supplies the numbers of \cref{sec:num-convergence}.
\lstinputlisting[language=Python]{code/analysis/S6_s6_part1_numerics.py}

\subsection{\texttt{S6\_s6\_part2\_robust.py}}
Part 2 of S6 produces the tables of fixed-current metrics for the final converged comparison, the isolation tests
B4--B7 (\run{0048}--\run{0051}) with a rigid-shift test, and the closure of the \SI{6.3}{nm} hypotheses A1--A7, B0 and
C1. Curves are rescaled by the exact linearity in $\muband$ (verified in part 1 to $10^{-5}$).
\lstinputlisting[language=Python]{code/analysis/S6_s6_part2_robust.py}

\subsection{\texttt{S6\_s6\_part3\_predictions.py}}
Part 3 of S6 performs the computations of the prediction register from real ATLAS outputs: $\Id$--$\Vd$ (B8--B10,
\run{0041}--\run{0043}) and elevated temperature (B11--B14, \run{0044}--\run{0047}). The \SI{300}{K} references B1--B3 are
rescaled exactly to the mobility used in the prediction runs, in the two temperature variants P-phonon (as simulated,
TMUN $= 1.5$) and P-MTR (temperature-independent band mobility).
\lstinputlisting[language=Python]{code/analysis/S6_s6_part3_predictions.py}

\subsection{\texttt{S6\_s6\_write\_register.py}}
The script writes \file{PREDICTIONS_REGISTER.md} and \file{predictions_canonical.csv} from the outputs of part 3,
together with hashes; it makes no simulator call and modifies no package file (\cref{ch:app-register}).
\lstinputlisting[language=Python]{code/analysis/S6_s6_write_register.py}

\subsection{\texttt{REVIEW\_review\_verify.py}}
The script contains the referee's re-computation of the load-bearing specialist claims of \file{REVIEW.md}. It reads only
raw files (measured data, run outputs and \file{deckbuild.out}) and imports \file{scripts/extract_metrics.py} unchanged;
no ATLAS launch is made.
\lstinputlisting[language=Python]{code/analysis/REVIEW_review_verify.py}

\subsection{\texttt{SYNTHESIS\_syn\_checks.py}}
The script contains the read-only checks of the synthesis stage. It is run from the package root with the Python
environment of the project, and its outputs are written to \file{scratch/SYNTHESIS/syn_checks_out.txt}.
\lstinputlisting[language=Python]{code/analysis/SYNTHESIS_syn_checks.py}

\subsection{\texttt{SYNTHESIS\_syn\_leverage.py}}
The script makes a read-only comparison of the leverage of the confinement package with that of the fitted and assumed
parameters at \SI{2}{nm} on the same metrics: fixed-current swings ($10^{-11}$--$10^{-10}$, $10^{-10}$--$10^{-9}$),
$\SScc$ ($10^{-10}$--$10^{-8}$) and the non-rigidity of the horizontal shift (shift at $10^{-8}$ minus shift at
$10^{-11}\Aum$), for the one-at-a-time pairs relative to \run{0012} or \run{0038}.
\lstinputlisting[language=Python]{code/analysis/SYNTHESIS_syn_leverage.py}

% -------------------------------------------------------------------
\section{Collection script of this report}\label{sec:code-collect}

The report itself is assembled from package files by scripts in \file{REPORT/tools/}; they read package files and write
only inside the report folder. In stage 0a, \file{tools/build_digest.py} builds \file{REPORT/digest/}, which contains the
paper texts split by printed page (\cref{sec:lit-convention}), the list of literature items used, the citations found in
the analysis reports, the ATLAS manual pages of the statements used by the decks, the run table \file{digest/runs.csv}
and a list of key numbers, so that writers never re-read raw sources. In stage 0b, \file{tools/make_figures.py} builds
every registry figure from real run and data files and writes \file{figures/<stem>.pdf}, a PNG copy and
\file{figures/MANIFEST.md} with the sources and key numbers of each figure; the call
\texttt{python tools/make\_figures.py [stem ...]} rebuilds selected figures. In stage 0c,
\file{tools/make_appendix_assets.py} copies code, decks, configuration and data into the report (ASCII-sanitised) and
generates \file{appendices/A_runs_table.tex} (\cref{tab:run-catalogue}) and \file{appendices/D_hashes_table.tex}
(re-verified hashes). Finally, \file{tools/lint_latex.py} is a static checker of the \LaTeX{} sources, since no TeX
installation is available on the project machine; it prints a report and exits with code 1 if any error is found.

\begin{keyresult}
All deck generation, solver launches, metric extraction and analysis are performed by the scripts listed here. Only three
package scripts launch ATLAS, and each launch is logged. All analysis scripts are read-only with respect to the runs and
share one metric extractor, so that measured and simulated metrics are computed identically. Surrogate models are
confined to decomposition and hypothesis generation (\evSURR).
\end{keyresult}
